% !TeX root = ../教材题库_全部条目_转写工作稿.tex
% Source: XJ-B1-C05-EX-P166-Y0573
% 第5章 / 5.2 任意角的三角函数 / 教材 p.166 / PDF 第173页
% OCR未校对：公式与文字均可能有误；未完成数学TeX转写；禁止直接发给学生
% 原题及相邻原解截图（核对用）：../题目原图/XJ-B1-C05-EX-P166-Y0573-01.webp
% 本文件只在转写完成后作为题干使用；OCR 条目尚不是可靠题干。
\begin{bankitem}{5.2 例题 例题3}{第5章 / 5.2 任意角的三角函数 / 教材 p.166 / PDF 第173页}{XJ-B1-C05-EX-P166-Y0573}{OCR未校对：公式与文字均可能有误；未完成数学TeX转写；禁止直接发给学生}
\begin{ocrdraft}
[教材p.166 / PDF第173页]
\par
0二0]设0<0且tan0>0，试确定0是第几象限的角．
\par
解因为sin0<0，
\par
所以0是第三或第四象限的角或终边在3，轴的负半轴上．
\par
因为tan0>0，
\par
所以0是第一或第三象限的角．
\par
因此，满足sin0<0且tan0>0的0是第三象限的角．
\end{ocrdraft}
\end{bankitem}
